% ===== PRÉAMBULE CANONIQUE — à coller VERBATIM en tête de CHAQUE .tex =====
\documentclass[11pt,a4paper]{article}
\usepackage[utf8]{inputenc}\usepackage[T1]{fontenc}\usepackage{lmodern}
\usepackage[french]{babel}
\usepackage{amsmath,amssymb,mathtools}\usepackage{icomma}
\usepackage[version=4]{mhchem}          % \ce{...} : équations & formules
\usepackage{siunitx}\sisetup{locale=FR} % \qty{}{}, \si{}, \ang{}
\usepackage{chemfig}                    % \chemfig{...} : structures développées
\usepackage{xcolor}\usepackage{colortbl}
\usepackage{tikz}\usepackage{pgfplots}\pgfplotsset{compat=1.18}
\usepackage[most]{tcolorbox}\usepackage{enumitem}\usepackage{array,booktabs}
\usepackage[a4paper,margin=2.2cm]{geometry}
\usetikzlibrary{arrows.meta,calc,angles,quotes,positioning,patterns,backgrounds,shapes.geometric}
\definecolor{primary}{RGB}{222,49,99}\definecolor{primarydark}{RGB}{150,22,55}
\definecolor{defcol}{RGB}{0,114,178}\definecolor{methcol}{RGB}{52,92,148}
\definecolor{excol}{RGB}{0,135,90}\definecolor{attncol}{RGB}{200,40,55}
\tcbset{boxbase/.style={enhanced,breakable,boxrule=0.9pt,arc=2.5pt,
  left=3mm,right=3mm,top=2.3mm,bottom=2.3mm,fonttitle=\bfseries,coltitle=white}}
% --- boîtes TUTORIEL (noms EXACTS, ne pas renommer) ---
\newtcolorbox{cle}[1][]{boxbase,colback=defcol!6,colframe=defcol,
  title={\textsc{À retenir}\ifx\relax#1\relax\relax\else\textmd{ : \itshape #1}\fi}}
\newtcolorbox{methode}[1][]{boxbase,colback=methcol!7,colframe=methcol,sharp corners,
  title={\textsc{Méthode}\ifx\relax#1\relax\relax\else\textmd{ --- \itshape #1}\fi}}
\newtcolorbox{exemplebox}[1][]{boxbase,colback=excol!5,colframe=excol,
  title={\textsc{Exemple}\ifx\relax#1\relax\relax\else\textmd{ : \itshape #1}\fi}}
\newtcolorbox{piege}[1][]{boxbase,colback=attncol!6,colframe=attncol,
  title={\textsc{Piège}\ifx\relax#1\relax\relax\else\textmd{ : \itshape #1}\fi}}
\newtcolorbox{remarque}[1][]{boxbase,colback=black!4,colframe=black!45,
  title={\textsc{Remarque}\ifx\relax#1\relax\relax\else\textmd{ : \itshape #1}\fi}}
% --- boîtes EXERCICE / CORRIGÉ (noms EXACTS, auto-comptées) ---
\newtcolorbox[auto counter]{exo}[1][]{boxbase,colback=primary!4,colframe=primary,
  title={\textsc{Exercice}~\thetcbcounter\ifx\relax#1\relax\relax\else\textmd{ --- \itshape #1}\fi}}
\newtcolorbox[auto counter]{corrige}[1][]{boxbase,colback=excol!4,colframe=excol,
  title={\textsc{Corrigé}~\thetcbcounter\ifx\relax#1\relax\relax\else\textmd{ --- \itshape #1}\fi}}
\newcommand{\R}{\mathbb{R}}\newcommand{\N}{\mathbb{N}}\newcommand{\Z}{\mathbb{Z}}\newcommand{\ds}{\displaystyle}
% (un \usepackage{fancyhdr}+en-tête est possible ; il est ignoré côté web)
% =========================================================================
\begin{document}
\begin{center}\color{primarydark}\large\bfseries Thème 3 — Potentiels standard et prévision \quad$\bullet$\quad Niveau Facile\end{center}

\begin{exo}[classer les oxydants]
On donne : $E_{\mathrm{r}}^{\circ}(\ce{Zn^{2+}}/\ce{Zn})=-0{,}76$~V ;
$E_{\mathrm{r}}^{\circ}(\ce{Cu^{2+}}/\ce{Cu})=+0{,}34$~V ;
$E_{\mathrm{r}}^{\circ}(\ce{Ag+}/\ce{Ag})=+0{,}80$~V ;
$E_{\mathrm{r}}^{\circ}(\ce{F2}/\ce{F^-})=+2{,}87$~V.
Classe les oxydants \ce{Zn^{2+}}, \ce{Cu^{2+}}, \ce{Ag+}, \ce{F2} du \textbf{plus faible} au
\textbf{plus fort}.
\end{exo}

\begin{exo}[le plus fort oxydant, le plus fort réducteur]
On donne : $E_{\mathrm{r}}^{\circ}(\ce{Al^{3+}}/\ce{Al})=-1{,}66$~V ;
$E_{\mathrm{r}}^{\circ}(\ce{Fe^{2+}}/\ce{Fe})=-0{,}44$~V ;
$E_{\mathrm{r}}^{\circ}(\ce{Cl2}/\ce{Cl^-})=+1{,}36$~V.
\begin{enumerate}[label=\alph*)]
  \item Quel est l'oxydant le plus fort ?
  \item Quel est le réducteur le plus fort ?
\end{enumerate}
\end{exo}

\begin{exo}[réaction spontanée ou non ?]
Données : $E_{\mathrm{r}}^{\circ}(\ce{Cu^{2+}}/\ce{Cu})=+0{,}34$~V ;
$E_{\mathrm{r}}^{\circ}(\ce{Zn^{2+}}/\ce{Zn})=-0{,}76$~V. Chaque réaction est-elle spontanée dans les
conditions standard ? Justifie par le signe de $\Delta E^{\circ}$.
\begin{enumerate}[label=\alph*)]
  \item \ce{Cu^{2+} + Zn -> Cu + Zn^{2+}}
  \item \ce{Zn^{2+} + Cu -> Zn + Cu^{2+}}
\end{enumerate}
\end{exo}

\begin{exo}[quel métal réagit avec un acide ?]
On plonge un métal dans l'acide chlorhydrique (couple \ce{H+}/\ce{H2}, $E_{\mathrm{r}}^{\circ}=0{,}00$~V).
Y a-t-il \textbf{dégagement de dihydrogène} ? Données : $E_{\mathrm{r}}^{\circ}(\ce{Zn^{2+}}/\ce{Zn})
=-0{,}76$ ; $E_{\mathrm{r}}^{\circ}(\ce{Fe^{2+}}/\ce{Fe})=-0{,}44$ ;
$E_{\mathrm{r}}^{\circ}(\ce{Cu^{2+}}/\ce{Cu})=+0{,}34$ ;
$E_{\mathrm{r}}^{\circ}(\ce{Ag+}/\ce{Ag})=+0{,}80$~V.
\begin{enumerate}[label=\alph*)]
  \item le zinc \ce{Zn}
  \item le fer \ce{Fe}
  \item le cuivre \ce{Cu}
  \item l'argent \ce{Ag}
\end{enumerate}
\end{exo}

\begin{exo}[écrire la réaction spontanée]
Deux couples : $E_{\mathrm{r}}^{\circ}(\ce{Ag+}/\ce{Ag})=+0{,}80$~V et
$E_{\mathrm{r}}^{\circ}(\ce{Cu^{2+}}/\ce{Cu})=+0{,}34$~V.
\begin{enumerate}[label=\alph*)]
  \item Lequel des deux couples fournit l'oxydant ?
  \item Écris l'équation de la réaction spontanée.
\end{enumerate}
\end{exo}

\end{document}
